%% ==========================================================================
%% SECTION 2: BACKGROUND
%% ==========================================================================
\section{Background: Why Truth Emerges (and Where It Breaks)}\label{sec:background}

The training corpus of a large language model is a fossil record of human
communication. Scientific papers report experimental results. Wikipedia articles
document verifiable claims. Legal filings assert facts under penalty of perjury.
The majority of this text was produced in institutional contexts where accuracy
was expected, if not always achieved.

This reflects mixed selection pressures on human communication. Coordination
demands truth-tracking: markets require reliable signals, contracts require
shared semantics, science requires reproducibility. But status competition and
social bonding reward fluent, engaging speech even when not fully accurate. The
training corpus contains both: institutional text that is broadly truth-tracking,
and misinformation, bias, and fabrication.

A language model trained on this corpus inherits these mixed regularities. Base
models hallucinate substantially on factual benchmarks because they are not honest by
default. But they also encode rich internal representations of concepts like
correctness and truthfulness, which can be elicited under the right conditions.
The base model's relationship to truth is complex: neither reliably honest nor
uniformly deceptive.

RLHF can strengthen or weaken this inheritance, depending on implementation.

Specialized designs like TruthRL, which use ternary rewards distinguishing
``correct,'' ``hallucinated,'' and ``abstain'', show measurable improvements:
28.9\% hallucination reduction and 21.1\% truthfulness improvement on
knowledge-intensive benchmarks~\cite{truthrl}. These results demonstrate that
preference optimization \emph{can} teach models to recognize their knowledge
boundaries when the reward signal is carefully constructed.

Public descriptions of major RLHF pipelines rarely describe such ternary rewards,
and available audits highlight underweighting of abstention and explicit
truthfulness signals.

Documented RLHF practice often conflates helpfulness, harmlessness, and honesty
under vague ``HHH'' criteria, with definitional slippage obscuring trade-offs
between user-friendliness and truthfulness. Human feedback is progressively
bootstrapped into AI-generated feedback, reducing direct supervision of factual
correctness. Current popular benchmarks and leaderboards give zero or minimal
credit for abstention, making confident guessing more score-optimal than
calibrated uncertainty.

Under these conditions, the optimization pressure favors fluent completion over
honest refusal. This can be formalized as a reward inequality:

\begin{multline*}
w_{\text{coherence}} R_{\text{coherence}} + w_{\text{engagement}} R_{\text{engagement}} \\
\gg w_{\text{factuality}} R_{\text{factuality}} + w_{\text{refusal}} R_{\text{refusal}}
\end{multline*}

No direct measurement of these weights in commercial reward models exists. But
current evaluation and feedback practices incentivize behavior consistent with
such an inequality. The False-Correction Loop emerges: models apologize, claim
correction, and fabricate again, because conversational fluency dominates
factual accuracy in the effective reward signal.

The problem is not that RLHF cannot prioritize honesty. Specialized designs
prove otherwise. The problem is that under realistic mixed-objective incentives
and bounded oversight, particularly in hard cases where ground truth is
unavailable and verification is expensive, supervisors operating under current
generic protocols struggle to reliably distinguish truthful uncertainty from
confident fabrication.

This is the gap our impossibility result formalizes.
